\chapter{REPOSITÓRIO DOS CÓDIGOS COMPUTACIONAIS}\label{ap:codigos}
\setcounter{figure}{0}\renewcommand{\thefigure}{A.\arabic{figure}}
\setcounter{table}{0}\renewcommand{\thetable}{A.\arabic{table}}

Os códigos desenvolvidos e adaptados para o processamento dos dados e para a reprodução dos cálculos deste trabalho estão disponíveis no repositório indicado a seguir. O Quadro~\ref{qua:resultado-A-1} apresenta os arquivos disponibilizados e suas principais funções. Os dados sísmicos, os registros de poços e os demais arquivos de entrada não são distribuídos no repositório em razão das condições de acesso e do volume dos dados.

Link: \url{https://github.com/mariaeulalia2002/codigos_tcc}

\begin{quadro}
  \centering
  \caption[Códigos disponibilizados no repositório e suas principais funções]{Códigos disponibilizados no repositório e suas principais funções.}
  \label{qua:resultado-A-1}
  \small
  \setlength{\tabcolsep}{3pt}
  \renewcommand{\arraystretch}{1.08}
  \begin{tabularx}{\linewidth}{|>{\raggedright\arraybackslash}p{0.15\linewidth}|>{\raggedright\arraybackslash}X|>{\raggedright\arraybackslash}p{0.20\linewidth}|>{\centering\arraybackslash}p{0.25\linewidth}|}
    \hline
    \textbf{\textbf{Código}} & \textbf{\textbf{Principal função}} \\
    \hline
    verificacao\_headers\_segy.py & Comparar o intervalo amostral registrado no cabeçalho binário com os valores presentes nos cabeçalhos dos traços dos arquivos SEG-Y. \\ \hline
    analise\_Espectral.py & Calcular os espectros de potência nos pontos de controle, determinar as frequências dominantes e fornecer os valores utilizados na avaliação da resolução vertical sísmica. \\ \hline
    extracao\_dados\_vsp\_poco\_137.py & Extrair do arquivo DLIS do poço 1-BRSA-137-RJS os dados de profundidade e tempo utilizados no cálculo das velocidades intervalares. \\ \hline
    suavizacao\_poco\_137.py & Aplicar média móvel simples, filtro Savitzky--Golay e média móvel ponderada à curva de velocidades do poço 1-BRSA-137-RJS. \\ \hline
  \end{tabularx}
\end{quadro}
\begin{quadro}
  \centering
  \begin{center}\footnotesize Quadro~\thequadro{} — Códigos disponibilizados no repositório e suas principais funções. (continuação).\end{center}\smallskip
  \small
  \setlength{\tabcolsep}{3pt}
  \renewcommand{\arraystretch}{1.08}
  \begin{tabularx}{\linewidth}{|>{\raggedright\arraybackslash}p{0.15\linewidth}|>{\raggedright\arraybackslash}X|>{\raggedright\arraybackslash}p{0.20\linewidth}|>{\centering\arraybackslash}p{0.25\linewidth}|}
    \hline
    \textbf{\textbf{Código}} & \textbf{\textbf{Principal função}} \\
    \hline
    conversao\_tempo\_profundidade.py & Calcular as velocidades equivalentes dos intervalos sedimentares e a ponderação das velocidades dos poços pelo inverso da distância aos eventos. \\ \hline
    identificacao\_picos\_secundarios.py & Identificar possíveis picos secundários f1 e f2 nas curvas HVSR previamente processadas pelo \textit{hvsrweb/hvsrpy}. \\ \hline
    incertezas\_morfometricas/incertezas\_morfometricas.py & Calcular os valores convertidos de espessura, soterramento, volume, altura total de queda e inclinação, além de propagar as respectivas componentes de incerteza. \\ \hline
  \end{tabularx}
  \nota{*\textit{ o processamento principal das curvas HVSR foi realizado com os }notebooks\textit{ do }hvsrweb\textit{ e com a biblioteca }hvsrpy*, de \citeonline{vantassel2021}, que não são redistribuídos neste repositório. A rotina de identificação de picos secundários corresponde ao procedimento complementar desenvolvido para este trabalho.}
\end{quadro}
